\par\noindent\textbf{(1)} 环所链磁通 \(\Phi_B(t)=\pi a^2B(t)\)。Faraday定律 \(\mathcal E=-d\Phi_B/dt=-\pi a^2\dot B(t)\)，负号体现Lenz反抗变化；在圆对称理想化中环上切向电场平均为 \(E_\phi=\mathcal E/(2\pi a)=-a\dot B/2\)。
\par\noindent\textbf{(2)} 环总电阻 \(R=2\pi ar_\ell\)。忽略自感时 \(I(t)=\mathcal E/R=-a\dot B/(2r_\ell)\)，焦耳功率 \(P=I^2R=\mathcal E^2/R=\pi a^3\dot B^2/(2r_\ell)\ge0\)。此结果用于外磁场缓慢驱动，不能忽略在高频产生的涡流自感。
\par\noindent\textbf{(3)} 环磁矩 \(m_z=I\pi a^2=-\pi a^3\dot B/(2r_\ell)\)，与磁场变化方向相反。外源施加磁通变化的电功 \(I\mathcal E=I^2R=P\)，因此能量以热的形式耗散，没有出现负功率。
\par\noindent\textbf{(4)} 若自感为 \(L>0\)，环路满足 \(L\dot I+RI=-\pi a^2\dot B\)。其电磁能量恒等式 \(d(LI^2/2)/dt+RI^2=I\mathcal E\) 显示一部分外源功存入磁能。只有驱动角频率 \(\omega L/R\ll1\) 才可退化到 \(I\approx\mathcal E/R\)。
